\section{Integral cycles and their attachment}\label{sec:integral-attachments}

All cohomology in this section has integral coefficients
unless otherwise specified. The Mukai marking identifies
the exceptional curve classes, but rational intersection
formulas alone do not determine how their neighborhoods
attach to the complement. We calculate both the primitive
integral classes and the resulting boundary homomorphisms.

\subsection{Primitive classes in the isolated example}
On $P_+$ the universal extension has $K$-class
$b\,\OO(1)+a$, up to tensoring by a base line bundle.
For $x\in v^\perp$ this ambiguity disappears from
\eqref{eq:mukai-marking}; hence
\[
 \theta_+(x)|_{P_+}=(a,x)h,\qquad
 \theta_-(x)|_{P_-}=-(a,x)h.
\]
We identify a rational vector in $v^\perp$ with its
curve functional by the Mukai pairing. A plane line
therefore has class
\begin{equation}\label{eq:isolated-classes}
 \ell_\pm=\pm(a-v/2),\qquad q(\ell_\pm)=-5/2,
 \qquad \rho_\pm=\theta_\pm\bigl(\pm(2a-v)\bigr).
\end{equation}
The integral vector $w_0=(9,3D,-10)$ lies in $v^\perp$
and satisfies $(a,w_0)=1$. It proves that each line
is primitive. The vector $2a-v=(-1,-D,6)$ is primitive;
its pairings with $v^\perp$ are even, and its pairing
with $w_0$ is two. Thus $q(\rho_\pm)=-10$ and
$\rho_\pm$ has divisibility two. The restriction of
$\theta_\pm(w_0)^2$ to the plane is $h^2$, which
also proves primitivity of its surface class.

For a K3$^{[2]}$-type fourfold the polarized Fujiki
identities are
\begin{equation}\label{eq:fujiki}
 \begin{aligned}
 \int_M\alpha\beta\gamma\delta
 &=q(\alpha,\beta)q(\gamma,\delta)
   +q(\alpha,\gamma)q(\beta,\delta)
   +q(\alpha,\delta)q(\beta,\gamma),\\
 \int_M c_2(M)\alpha\beta&=30q(\alpha,\beta).
 \end{aligned}
\end{equation}
The rational map $\operatorname{Sym}^2H^2(M,\mathbb Q)
\to H^4(M,\mathbb Q)$ is an isomorphism. Since
$q(\ell)=-5/2$, the class $\ell^2/2+c_2(M)/24$
pairs with $\alpha\beta$ as
$q(\ell,\alpha)q(\ell,\beta)$, exactly the
restriction pairing on the plane. Nondegeneracy gives
the classical formula
\begin{equation}\label{eq:plane-class}
 [P_\pm]=\frac{\ell_\pm^2}{2}+\frac{c_2(M_\pm)}{24}
         =\frac{\rho_\pm^2}{8}+\frac{c_2(M_\pm)}{24}.
\end{equation}
See \cite{hassetttschinkel}, with the attribution
there to O'Grady. The cotangent normal bundle gives
$[P]^2=3$; from $T_M|_P=T_P+T^*_P$ in $K$-theory
one also gets $\int_Pc_2(M)=-3$. Formula
\eqref{eq:plane-class} is rational. The integral
primitivity assertion follows from $w_0$, without
assuming an integral symmetric-square presentation.

Totaro's torsion-freeness theorem
\cite[Theorem~3.1]{totaro} and the coefficient of
degree two in G\"ottsche's Hilbert-scheme formula
\cite{gottsche} give free even cohomology of ranks
$1,23,276,23,1$ for $M_\pm$. The contraction identifies
$Z$ homeomorphically with $M_\pm/P_\pm$: the map
from this compact quotient to the Hausdorff target
is bijective. The plane inclusion is a cofibration,
and its restrictions in degrees 0, 2 and 4 are
surjective. The pair sequence therefore gives
\begin{center}
\begin{adjustbox}{max width=\textwidth}
\begin{tabular}{lrrrrr}
\toprule
Space & $b_0$ & $b_2$ & $b_4$ & $b_6$ & $b_8$\\
\midrule
$M_\pm$ & 1&23&276&23&1\\
$Z$ &1&22&275&23&1\\
$M_\pm\setminus P_\pm$ &1&23&275&22&0\\
\bottomrule
\end{tabular}
\end{adjustbox}
\end{center}
All other groups vanish, and these groups are
torsion-free. The pullback from $Z$ is an injective
ring map whose image is
\begin{equation}\label{eq:isolated-target-ring}
 \Z\cdot1\ \oplus\left(\bigoplus_{k>0}
 \ker\bigl(H^k(M_\pm)\longrightarrow H^k(P_\pm)\bigr)\right).
\end{equation}
Its multiplication is inherited from the chamber
ring; the local singularity alone does not provide
that ring as additional output.

\subsection{The two interacting curve and surface classes}
Return to the chamber $\zeta_0,\zeta_2>0$, and put
$E=P_0\cup P_2$. Let $c_i$ be an effective line in
$P_i$ and $f$ a moving projective-line fibre. The
universal extensions on these fibres give
\begin{equation}\label{eq:interacting-curve-classes}
 \begin{gathered}
 l_i=a_i-v/2\quad(i=0,2),\qquad
 c_0=-l_0,\quad c_2=-l_2,\quad f=a_1,\\
 f=c_0+c_2,\qquad
 q(c_0)=q(c_2)=-5/2,\quad q(c_0,c_2)=3/2,\quad q(f)=-2.
 \end{gathered}
\end{equation}
Indeed the restriction to $P_i$ is
$\theta_v(x)|_{P_i}=-(a_i,x)h_i$. Over a moving
fibre the extension is $F\,\OO(1)+L_1$, so its
restriction is $-(w,x)h=(a_1,x)h$. This calculation
uses a fixed $F$ and needs no global ordinary
universal family on $S^\circ$. Equation
\eqref{eq:reduced-specialization} realizes the class
relation by a reduced family of actual curves.

Use the primitive lattice embedding
\eqref{eq:primitive-embedding}. The integral vectors
\[
 x=(0,f_2-f_3,0),\qquad y=(0,-3f_2,-1)
\]
are in $v^\perp$, with pairings $(0,1,-1)$ and
$(1,-2,1)$ against $(a_0,a_1,a_2)$.
The same witnesses exist for every primitive embedding of
$\NS(X)$: unimodularity of $H^2(X,\Z)$ makes the restriction
$H^2(X,\Z)\to\NS(X)^*$ surjective, so the prescribed
pairings of their degree-two components can always be realized.
For
$\alpha=\theta_v(x)$ and $\beta=\theta_v(y)$,
\begin{equation}\label{eq:restriction-witnesses}
 \alpha|_E=(0,h_2),\qquad \beta|_E=(-h_0,-h_2),
 \qquad
 \alpha^2|_E=(0,h_2^2),\quad
 \beta^2|_E=(h_0^2,h_2^2).
\end{equation}
Both restriction matrices are unimodular. Thus the
two line classes span a primitive subgroup of
$H_2(M)$, and the two plane classes span a primitive
subgroup of $H_4(M)$.

The cotangent normal bundles and the transverse
intersection at $p$ give
\begin{equation}\label{eq:two-plane-matrix}
 B=\begin{pmatrix}3&1\\1&3\end{pmatrix}.
\end{equation}
Equivalently, $[P_i]=l_i^2/2+c_2(M)/24$ by the
argument for \eqref{eq:plane-class}; the off-diagonal
intersection is $q(l_0,l_2)^2/2-1/8=1$.
Here the geometric intersection and the rational
Fujiki calculation agree, while
\eqref{eq:restriction-witnesses} supplies the integral
information.

\subsection{The neighborhood boundaries}
For one plane the boundary is the sphere bundle
$L=S(T^*\PP^2)$. Its Euler class is $3h^2$.
The Gysin sequence gives $H^j(L)=\Z$ for
$j=0,2,5,7$, $H^4(L)=\Z/3$, and zero otherwise.
In particular $H^3(L)=0$, although $H_3(L)=\Z/3$.

For two transverse planes take a closed regular
neighborhood $W$ with smoothed boundary $L$. It
is the plumbing of their normal disk bundles and
retracts to $E\simeq\PP^2\vee\PP^2$.
Poincar\'e--Lefschetz duality makes the middle map
$H^4(W,L)\to H^4(W)$ the matrix $B$.
The pair sequence, together with $\det B=8$ and
the gcd one of its entries, gives
\begin{equation}\label{eq:two-plane-boundary}
 H^j(L)=
 \begin{cases}
 \Z,&j=0,7,\\
 \Z^2,&j=2,5,\\
 \Z/8,&j=4,\\
 0,&\text{otherwise}.
 \end{cases}
\end{equation}
The quotient map on $H^4(W)=\Z^2$ is
\begin{equation}\label{eq:boundary-functional}
 (n_0,n_2)\longmapsto n_0-3n_2\pmod8.
\end{equation}
Its kernel is $B\Z^2$: it kills both columns and
is onto a group of order $\det B$.
This $L$ is a smooth boundary in the resolution.
The target has a positive-dimensional singular
stratum through $o$, so $L$ is not the link of an
isolated target singularity.

Alexander--Lefschetz duality gives
$H^k(M,M\setminus E)=H_{8-k}(E)$. The primitive
line, plane and point inclusions just proved show
that $V=M\setminus E$ has torsion-free even
cohomology, with ranks $1,23,274,21$ in degrees
$0,2,4,6$. In particular,
\begin{equation}\label{eq:complement-quotient}
 H^4(V)=H^4(M)/(\Z[P_0]+\Z[P_2]).
\end{equation}
Notice that $V$ retains the moving exceptional
divisor away from $E$.

\subsection{Calculated attachments on integral generators}
We now specify the map from the complement to the
boundary. The input includes the complement ring
with its distinguished universal-index classes.
The following formulas evaluate the attachment
without asking for an ambient lift of a given class.

For $z\in K^0_{\mathrm{top}}(X)$ define the family
index
\begin{equation}\label{eq:universal-index}
 I_z=p_{M!}\bigl(p_X^*(z^\vee)\otimes\cE\bigr)
       \in K^0_{\mathrm{top}}(M),\qquad
 \chi(F,z)=-(v(F),v(z)).
\end{equation}
Here $p_!$ is the complex $K$-theory pushforward.
Markman's integral-generation theorem
\cite[Theorem~1]{markman}, the integral K\"unneth
decomposition and the perfect Euler pairing on
$K^0_{\mathrm{top}}(X)$ show that the Chern classes
of these indices generate $H^*(M,\Z)$. Thus degree
four is generated by the classes $c_2(I_z)$ and
$c_1(I_z)c_1(I_{z'})$. Rational Chern characters
would not suffice for this conclusion.

One normalization removes the base-line ambiguity
of the universal family on the exceptional set.
For the isolated example, let $F,Q$ be its
subobject and quotient in the chosen chamber.
Their Mukai vectors span a primitive rank-two
sublattice: the rank-one factor and the primitive
$D$-coordinate give a unit minor. Unimodularity
supplies a topological class $y$ with
$\chi(F,y)=0$, $\chi(Q,y)=1$. On the plane
\[
 \cE=(F\,\OO(1)+Q)\otimes L,
 \qquad I_y|_P=L.
\]
The global index $I_y$ has rank one. Tensor
$\cE$ by $\det(I_y)^{-1}$ to remove $L$.
For this normalized topological universal class,
\begin{equation}\label{eq:isolated-index-restriction}
 c(I_z)|_P=(1+h)^{m_z},\qquad m_z=\chi(F,z).
\end{equation}
The powers are interpreted as total Chern classes
of virtual bundles and allow negative integers.

For the interacting planes the class
$y=(0,-f_2,0)$ has Euler pairings $(0,1,0)$ with
$(a_0,a_1,a_2)$. The same global determinant
normalization removes the line twists on \emph{both}
planes, including their common point. It gives
\begin{equation}\label{eq:interacting-index-restriction}
 c(I_z)|_{P_i}=(1+h_i)^{m_{i,z}},\qquad
 m_{i,z}=\chi(L_i,z),\quad i=0,2.
\end{equation}
This normalization is topological; it makes no
claim of an algebraic lift for the auxiliary
Brauer class considered in the next section.

\begin{proposition}\label{prop:calculated-attachments}
Let $A=H^*(V,\Z)$ be the complement ring with the
restrictions of the normalized index Chern classes
just specified.

For an isolated plane, the actual restriction
$\kappa:A^4\to H^4(L)=\Z/3$ is determined by
\begin{equation}\label{eq:mod-three-generators}
 \begin{aligned}
 \kappa(c_2(I_z)|_V)&=\binom{m_z}{2},\\
 \kappa(c_1(I_z)c_1(I_{z'})|_V)&=m_zm_{z'}
 \end{aligned}\qquad\pmod3,
\end{equation}
and restriction identifies
\begin{equation}\label{eq:mod-three-gluing}
 H^4(M)=\{(u,N)\in A^4\oplus\Z:
                      N\equiv\kappa(u)\pmod3\}.
\end{equation}
For the two-plane chamber, the actual restriction
$\kappa:A^4\to H^4(L)=\Z/8$ is determined by
\begin{equation}\label{eq:mod-eight-generators}
 \begin{aligned}
 \kappa(c_2(I_z)|_V)
   &=\binom{m_{0,z}}2-3\binom{m_{2,z}}2,\\
 \kappa(c_1(I_z)c_1(I_{z'})|_V)
   &=m_{0,z}m_{0,z'}-3m_{2,z}m_{2,z'}
 \end{aligned}\qquad\pmod8,
\end{equation}
and restriction identifies
\begin{equation}\label{eq:mod-eight-gluing}
 H^4(M)=\{(u,n_0,n_2)\in A^4\oplus\Z^2:
           n_0-3n_2\equiv\kappa(u)\pmod8\}.
\end{equation}
Both attachment maps are surjective.
\end{proposition}

\begin{proof}
The relative group $H^5(M,V)=H_3(E)$ vanishes
in both cases. Thus $H^4(M)\to A^4$ is onto,
and the restricted degree-four index monomials
generate $A^4$. Equations
\eqref{eq:isolated-index-restriction} and
\eqref{eq:interacting-index-restriction} give their
integer evaluations on the plane or planes.
Passing to the boundary gives exactly
\eqref{eq:mod-three-generators} and
\eqref{eq:mod-eight-generators}.

To check every relation among these generators,
let a degree-four integral polynomial in the
restricted Chern classes vanish on $V$. Its
ambient polynomial class lies in the relative
image $H^4(M,V)$, generated by the plane classes.
Its evaluations on $E$ therefore lie in
$3\Z$ or $B\Z^2$, respectively. The prescribed
boundary functional kills these subgroups.
Consequently the formulas are independent of the
chosen polynomial expression in $A^4$ and
define its actual boundary restriction.

Apply Mayer--Vietoris to $V$ and an open regular
neighborhood of $E$. Their overlap retracts to
$L$, and $H^3(L)=0$. The restriction from $M$
is therefore injective in degree four, with image
the kernel of the difference of the two boundary
maps. This is precisely the displayed fibre
product. Finally $\theta(w_0)^2$ evaluates to one
in the isolated case, while $\alpha^2$ evaluates
to $(0,1)$ in the interacting case. Their images
are units modulo three and eight.
\end{proof}

The construction retains the complement ring and
its marked index classes as input. It recovers
the integral degree-four attachment, including
the torsion boundary information that is absent
from Betti numbers. A change of index basis or
universal normalization changes the formulas for
the marked generators, but not the restriction
map they calculate.
